\begin{frame}{자주 묻는 질문 (1): ε를 크게?}
  \begin{block}{Q: $\varepsilon$ 을 크게 잡을수록 항상 더 좋은가?}
    \textbf{아니다.}
    \begin{itemize}
      \item $\varepsilon$ 이 0이면 최적해 영영 놓칠 위험.
      \item 지나치게 크면(예: 0.5) 이미 최선의 팔을 찾은 뒤에도
        나쁜 팔을 절반 확률로 계속 시도.
      \item 실전: 학습 진행될수록 $\varepsilon$ 을 점점 줄이는
        \textbf{decay} 가 널리 쓰이는 절절점.
    \end{itemize}
  \end{block}
  \vskip0.5em
  \begin{block}{Q: $\frac{1}{N_t(a)}$ 대신 고정 $\alpha$ 를 쓰면?}
    \begin{itemize}
      \item $\frac{1}{N}$: 시도 횟수 ↑ $\to$ 갱신 폭 ↓ $\to$
        정확한 평균으로 \textbf{수렴}.
      \item 고정 $\alpha$: 정확히 수렴하지 않으나,
        \textbf{최근 관찰}에 더 큰 가중치를 계속 유지.
        \textbf{비정상(non-stationary)} 환경에서 오히려 유리.
    \end{itemize}
  \end{block}
\end{frame}
